\chapter{Introduction}

\section{Background}
In recent years, millions of people in India have started investing in the stock market for the first time. In the past, buying shares was complicated and required physical paperwork. Today, anyone with a smartphone can open an investment account within minutes. With digital identity verification through Aadhaar and PAN (known as Know Your Customer, or KYC) and instant online payments using the Unified Payments Interface (UPI), investing has become accessible to ordinary citizens. 

According to official records published by the Securities and Exchange Board of India (SEBI), the number of electronic share-holding accounts (known as Demat accounts) in India increased from approximately 36 million in 2019 to more than 160 million by early 2026. A \textbf{Demat account} (short for Dematerialized account) works like a digital bank account, but instead of holding money, it holds electronic certificates of company shares and government bonds.

However, while opening an account is now easy, investing wisely remains difficult. Most new retail investors do not have formal financial training. When people start investing on their own, they are often overwhelmed by large amounts of confusing information:
\begin{itemize}
    \item Complex quarterly financial statements published by companies.
    \item Economic policy updates and interest rate announcements from the Reserve Bank of India (RBI).
    \item Thousands of news articles, social media tips, and daily stock price movements.
\end{itemize}

Without proper guidance, individual investors often make predictable mistakes. They buy stocks that have already risen sharply because of popular hype, sell good stocks in a panic during normal market declines, or put all their savings into one or two risky companies. In finance, this risk is called \textbf{lack of diversification}. Diversification simply means spreading investments across different companies and industry sectors so that poor performance in one stock does not wipe out an investor's entire savings.

To solve this problem, retail investors usually turn to mutual funds or automated investment tools called \textbf{robo-advisors}. A mutual fund pools money from many investors and pays a professional manager to select stocks. However, mutual funds charge recurring management fees (called \textbf{expense ratios}) of 1.5\% to 2.5\% every year. Over ten or twenty years, these fees take away a significant portion of an investor's accumulated returns. On the other hand, existing robo-advisors act like closed ``black boxes.'' They ask a few generic questions and suggest a standard basket of stocks or funds, without explaining why each company was chosen or how the investment plan connects to the user's specific life goals. When stock prices drop, investors do not understand why they hold these stocks, get frightened, and sell at a loss.

\section{Problem Statement}
To address these challenges, the problem statement of this Major Project is defined as follows:

\begin{quote}
\textit{``To design and develop an intelligent portfolio creation application that helps retail investors build personalized investment portfolios by providing transparent, context-aware, explainable, and mathematically reliable recommendations.''}
\end{quote}

\section{Challenges in Retail Portfolio Construction}
Building an automated application that creates personalized stock portfolios for Indian retail investors involves solving six key challenges:

\begin{enumerate}
    \item \textbf{Scattered and Overwhelming Information}: Financial information in India is spread across many different places: official stock exchange websites, company annual reports, news agencies, and government portals. An individual investor cannot read and make sense of hundreds of pages of filings every week.
    \item \textbf{Emotional Decision-Making}: Human investors are strongly influenced by emotion. Studies in behavioral finance show that people feel the pain of losing money twice as strongly as they enjoy making the same amount of money \cite{kahneman1979}. When markets fall, fear causes investors to sell at the worst possible time. An automated system must remain objective and rational.
    \item \textbf{The ``Black-Box'' Problem}: Existing automated advisory tools calculate investment percentages using hidden algorithms. They rarely explain why a particular stock was selected or why another was rejected. Without clear explanations, everyday investors find it difficult to trust automated advice.
    \item \textbf{Mathematical Mistakes in Generative AI}: Modern conversational Artificial Intelligence tools (Large Language Models, or LLMs) are good at reading articles and writing text. However, studies show that language models often make basic calculation mistakes when trying to calculate portfolio numbers \cite{kulpa2025}. For example, an LLM might suggest percentages that do not add up to 100\%, recommend negative percentages, or change its numbers unpredictably when asked the same question twice.
    \item \textbf{Indian Exchange Rule of Whole Shares}: In several foreign countries, investors can buy fractional shares (for example, 0.25 shares of a costly company). In India, exchange regulations strictly forbid fractional shares; an investor must buy complete, whole shares (such as 1, 2, or 10 shares). A stock like MRF Ltd. costs over Rs.~1,00,000 per share. If an algorithm suggests buying 1.5 shares, that order cannot be executed. The system must convert percentage targets into exact whole shares while ensuring that unallocated cash is safely kept in the investor's balance.
    \item \textbf{Managing a Large Market Universe}: The Indian equity market includes more than 2,600 actively traded companies on the National Stock Exchange of India (NSE) and the Bombay Stock Exchange (BSE). Evaluating thousands of companies for trading activity, financial stability, and price freshness requires fast and reliable data handling.
\end{enumerate}

\section{Motivation}
The primary motivation behind this project is to build an application called \textbf{Lucid}. The goal of Lucid is to give everyday retail investors the same quality of research, risk management, and mathematical precision that professional fund managers use, but in a simple, affordable, and transparent format.

Recent progress in \textbf{Agentic Artificial Intelligence} makes this possible \cite{hajaghaie2025}. Instead of using a single AI model to do everything, an agentic system divides the work among several specialized digital assistants (called \textbf{AI agents}). Each agent has a specific job:
\begin{itemize}
    \item One agent reviews broad economic conditions (such as inflation and interest rates).
    \item Another agent evaluates company balance sheets and profit trends.
    \item A third agent checks recent news and corporate filings.
    \item A risk critic agent double-checks the proposals to spot hidden risks.
\end{itemize}

Most importantly, Lucid keeps artificial intelligence and mathematical calculations strictly separate. The AI agents are used only for reading documents, evaluating qualitative trends, and shortlisting candidate companies. When it is time to calculate exact investment amounts and validate portfolio constraints, Lucid passes the task to a proven mathematical calculation engine. This separation ensures that the recommendations are both easy to understand and mathematically guaranteed to be accurate \cite{fatouros2024}.

\section{Project Objectives}
The specific technical and research objectives of this Major Project are:

\begin{enumerate}
    \item To develop an intelligent portfolio creation system that generates personalized investment recommendations based on an investor's financial goals, risk preferences, investment horizon, and other relevant factors.
    \item To ensure that the generated portfolio recommendations are reliable, transparent, explainable, and mathematically sound by considering relevant investor, market, and historical information.
    \item To validate and deploy the proposed system as a production-ready web application that provides an accessible interface for retail investors to create, explore, and understand their portfolios.
\end{enumerate}

\section{Scope of the Project}
The operational scope of Lucid defines what the application does and what it deliberately avoids:
\begin{itemize}
    \item \textbf{Market and Securities}: Lucid works with equity shares listed on the National Stock Exchange of India (NSE) and the Bombay Stock Exchange (BSE) in the standard cash delivery segment (where investors buy and hold real shares).
    \item \textbf{Target Users}: Individual Indian retail investors who want personalized portfolios tailored to their savings goals, health reviews of their existing stock holdings, and automated portfolio monitoring.
    \item \textbf{Order Quantities}: All recommended stock purchases are calculated as discrete whole shares (such as 5 shares or 12 shares) because fractional shares are not permitted on Indian exchanges.
    \item \textbf{Advisory Nature}: Lucid is an educational decision-support system. It explains its financial reasoning transparently so users can make informed choices.
    \item \textbf{Exclusions}: Lucid does not support intraday margin trading (borrowing money to buy and sell on the same day), high-risk derivative contracts (such as Futures and Options), or short selling (betting on a stock to fall). It also does not execute live broker trades automatically; it provides simulated trade previews and integration adapters for safe user testing.
\end{itemize}

\section{Proposed Solution: Lucid}
Lucid is built as a modular, web-based software platform designed from the ground up for clarity and ease of use:
\begin{itemize}
    \item \textbf{User Interface}: Built with modern web technologies (Next.js and React), providing clean screens where investors can enter their goals, view their portfolios, and ask questions.
    \item \textbf{Backend Server}: A high-speed application server (FastAPI) that handles user requests, coordinates security, and runs the application logic.
    \item \textbf{Multi-Agent Brain}: A team of specialized AI agents coordinated by an agent management system (LangGraph) that studies company filings, news, and economic reports.
    \item \textbf{Mathematical Engine}: A numerical solver (SciPy SLSQP) that calculates exact investment weights based on Modern Portfolio Theory, converts weights into whole share quantities, and accounts for Indian taxes and fees (such as Securities Transaction Tax and GST).
    \item \textbf{Continuous Monitoring Daemon}: A background program that quietly checks the investor's active portfolio during stock market hours (09:15 to 15:30 IST) to detect if stock price movements have caused the portfolio to drift away from its planned targets.
\end{itemize}

Lucid monitors the full universe of 2,639 listed Indian companies and filters them for trading liquidity and financial health, ensuring that every user receives recommendations tailored to their unique financial goals.

\section{Key Contributions}
The key contributions of this project are:
\begin{enumerate}
    \item \textbf{Separation of AI Reasoning from Math Calculations}: AI models write explanations and shortlist companies, while a mathematical solver calculates exact portfolio percentages, eliminating mathematical errors.
    \item \textbf{Whole-Share Integer Allocation}: An algorithm that calculates exact full-share numbers and protects unallocated cash balances, ensuring that every recommendation can be executed on Indian exchanges.
    \item \textbf{Automated Market-Hours Drift Tracking}: A monitoring service that checks portfolios every 5 minutes during trading hours to alert investors if stock weights shift significantly.
    \item \textbf{Clear, Plain-Language Explanations}: A feature called ``Stock-Why'' that explains in simple English exactly why each company was selected, alongside deeper financial metrics for advanced users.
    \item \textbf{Complete Indian Market Coverage}: A database covering all 2,639 listed companies on the NSE and BSE, rather than limiting users to a small, pre-selected list.
\end{enumerate}

\section{Report Organization}
The remainder of this report is organized as follows:
\begin{itemize}
    \item \textbf{Chapter 2 (Literature Review and Market Survey)} reviews existing financial theories, automated investment tools, and competing applications in India.
    \item \textbf{Chapter 3 (High-Level Design)} explains the overall structure of Lucid, how the major components fit together, and how an investor interacts with the system.
    \item \textbf{Chapter 4 (Low-Level Design)} details how the individual software modules and AI agent workflows operate internally.
    \item \textbf{Chapter 5 (Database Design)} describes the 62 database tables that store user information, stock data, and portfolio records.
    \item \textbf{Chapter 6 (System Architecture)} presents the eight software layers of Lucid, the technology stack, and runtime execution.
    \item \textbf{Chapter 7 (Methodology)} explains the step-by-step investment process, including mathematical formulas for portfolio optimization, lot sizing, and rebalancing.
    \item \textbf{Chapter 8 (Datasets Used)} explains what financial data Lucid uses, where it comes from, and how it is cleaned.
    \item \textbf{Chapter 9 (API Contract)} specifies how the web interface communicates with the backend server through structured requests and responses.
    \item \textbf{Chapter 10 (Work Done)} documents each implemented functional requirement from the user's perspective.
    \item \textbf{Chapter 11 (Results and Testing)} presents the testing methodology, system evaluation, and backtesting results.
    \item \textbf{Chapter 12 (Conclusions and Future Roadmap)} summarizes the achievements of Lucid and outlines future enhancements.
\end{itemize}
